\subsection{Introduction}

Much recent research has been dedicated to developing multi-modal drones that can transition between flight
and a stationary mode attached to a structure. Often referred to as "perching", this behavior is intended
to enable the use of drones over longer mission times, where perching can reduce power use, or in scenarios
where the noise of a drone's running motors is undesirable. Often the targeted use case for this is ecology research,
where drones could place a sensor payload in an inaccessible forest canopy, or monitor wildlife without the disruptive
noise of its motors \cite{liu_mini-drone_2023,meng_aerial_2022}. 



\subsection{Purpose}

This thesis aims to design a novel perch detection method that can detect cylindrical perches such as branches or power lines, 
as well as a grasping device that can attach a drone underneath such an object. The device should use low-cost components and accessible
software libraries to keep the proposed design affordable and accessible. This thesis will verify the reliability and accuracy
of the proposed detection method and the strength, and weight of the grasping mechanism. 


\subsection{Scope}

The scope of this thesis is limited to designing and verifying a perch detection and grasping device under controlled 
lab conditions. The device will not be integrated with a drone and no drone will be used in testing.


\subsection{Assumptions}

This thesis assumes that the device may need to support a drone with a mass of up to 500g, plus 31.58 N
of dynamic load. The dynamic load is to allow for force introduced by motion of the perch, such as a swaying tree branch.
Assuming that the drone is perched in a tree and subject to sinusoidal motion, its maximum acceleration is estimated at $63.17 m/s^2$ 
based on data provided by Kamimura et al. (0.4 m max displacement with a frequency of 2 Hz) \cite{kamimura_energy_2024}. 


\subsection{Significance}

Perching can increase the versatility of drones in data collection and inspection tasks, allowing them to achieve longer mission times
and deploy sensors in otherwise inaccessible locations. Novel and affordable detection ad grasping methods such as those proposed in this thesis can
make perching more accessible for a wider variety of applications.